% Contingut generat per llibre/llibre.py a partir de docs/equacions/.
% No l'editis a mà: torna'l a generar.

\chapter{Equacions i inequacions}

En aquest tema posem en pràctica tot el que hem après sobre nombres, potències, arrels i logaritmes per resoldre \textbf{equacions}: igualtats on hem de trobar els valors que les fan certes. Repassarem les equacions polinòmiques (de primer grau, segon grau, biquadrades i de grau superior), les racionals i irracionals, i les exponencials i logarítmiques. Després resoldrem \textbf{sistemes d'equacions} no lineals i \textbf{inequacions} de primer i de segon grau, on les solucions ja no són valors aïllats sinó intervals.

El mètode de resolució canvia segons el tipus d'equació, però la idea de fons és sempre la mateixa: transformar l'equació en una forma que ja sabem resoldre, i comprovar sempre les solucions obtingudes a l'equació original.

\section{Equacions polinòmiques}

\subsection{Definicions i conceptes bàsics}

Segurament ja has resolt moltes equacions abans. Comencem posant nom, de manera formal, al que és una equació i a la seva solució.

\begin{definicio}{Definició: equació i incògnita}
Sigui $f(x)$ una expressió algebraica. Una \textbf{equació} amb una incògnita és una igualtat que, amb les transformacions que calguin, es pot escriure de la forma

\[
f(x)=0.
\]

$x$ s'anomena la \textbf{incògnita} de l'equació.
\end{definicio}

\begin{exemple}{\textbf{Exemple:} Equacions}
Són exemples d'equacions:

\[
x-3=0, \qquad 5x-5=x+7.
\]

En el segon cas, la podem transformar:

\[
5x-5-x-7=0 \quad \Rightarrow \quad 4x-12=0,
\]

de manera que obtenim una expressió com la de la definició anterior.
\end{exemple}

\begin{definicio}{Definició: solució i equivalència}
Direm que un valor real $a$ és una \textbf{solució} de l'equació $f(x)=0$ si es compleix que

\[
f(a)=0.
\]

Si dues equacions tenen les mateixes solucions, direm que són \textbf{equivalents}.
\end{definicio}

\begin{exemple}{\textbf{Exemple:} Comprovació de solucions}
Si prenem les equacions de l'exemple anterior, en tots els casos $a=3$ n'és una solució:

\[
x-3=0 \;\longrightarrow\; 3-3=0,
\]

\[
5x-5=x+7 \;\longrightarrow\; 5\cdot 3-5=10=3+7,
\]

\[
4\cdot 3-12=0.
\]

Com que en tots els casos es tracta d'equacions de primer grau (que tenen una única solució), podem assegurar que les tres són equivalents.
\end{exemple}

\subsection{Equacions de primer i segon grau}

\begin{definicio}{Definició: equació polinòmica}
Una equació $p(x)=0$ és \textbf{polinòmica} si $p(x)$ és un polinomi. Parlarem aleshores del \textbf{grau de l'equació}, que coincideix amb el grau del polinomi.
\end{definicio}

\begin{exemple}{\textbf{Exemple:} Equacions de primer grau}
Una equació de primer grau es pot escriure de forma equivalent com

\[
ax+b=0, \qquad a\neq 0.
\]

Té una única solució, que es pot escriure de manera general com

\[
x=-\frac{b}{a}.
\]

Fixa't en el cas especial en què $a=0$ després de simplificar:

\[
2x-3=2(x-1) \;\Rightarrow\; 2x-3=2x-2 \;\Rightarrow\; -1=0.
\]

No existeix cap valor real que compleixi aquesta igualtat: l'equació no té solució.
\end{exemple}

\begin{exemple}{\textbf{Exemple:} Equacions de segon grau}
Una equació de segon grau es pot escriure com

\[
ax^2+bx+c=0, \qquad a\neq 0.
\]

Les solucions venen donades per la fórmula

\[
x=\frac{-b\pm\sqrt{b^2-4ac}}{2a},
\]

on $\Delta=b^2-4ac$ s'anomena \textbf{discriminant}. Segons el seu signe distingim tres casos:

\begin{itemize}
\item $\Delta>0$: dues solucions reals diferents.
\item $\Delta=0$: una solució real (doble).
\item $\Delta<0$: cap solució real.
\end{itemize}

Per exemple, a $4x^2+4x+1=0$ el discriminant és $\Delta=4^2-4\cdot4\cdot1=0$:

\[
x=\frac{-4\pm\sqrt{0}}{8}=-\frac{4}{8}=-\frac{1}{2} \quad \text{(solució doble)}.
\]
\end{exemple}

\subsection{Equacions de segon grau incompletes}

Una equació de segon grau $ax^2+bx+c=0$ és \textbf{incompleta} quan li falta el terme de primer grau ($b=0$) o el terme independent ($c=0$). En aquests casos no cal la fórmula: les podem resoldre de manera més ràpida.

\begin{definicio}{Cas 1: falta el terme de primer grau (b = 0)}
L'equació és de la forma $ax^2+c=0$. Aïllem $x^2$ i fem l'arrel quadrada, recordant que hi ha dos signes:

\[
ax^2+c=0 \;\Rightarrow\; x^2=-\frac{c}{a} \;\Rightarrow\; x=\pm\sqrt{-\frac{c}{a}}.
\]

Si $-\dfrac{c}{a}>0$, hi ha dues solucions oposades. Si $-\dfrac{c}{a}<0$, l'equació no té solució, perquè un quadrat no pot ser negatiu.
\end{definicio}

\begin{exemple}{\textbf{Exemple:} Equació incompleta sense terme de primer grau}
Resolem $2x^2-18=0$:

\[
2x^2=18 \;\Rightarrow\; x^2=9 \;\Rightarrow\; x=\pm\sqrt{9}=\pm 3.
\]

Les solucions són $x_1=3$ i $x_2=-3$.

En canvi, l'equació $x^2+4=0$ no té solució: ens quedaria $x^2=-4$, i cap nombre real elevat al quadrat dona un resultat negatiu.
\end{exemple}

\begin{definicio}{Cas 2: falta el terme independent (c = 0)}
L'equació és de la forma $ax^2+bx=0$. Traiem factor comú $x$ i fem servir que un producte val zero si i només si algun dels factors val zero:

\[
ax^2+bx=0 \;\Rightarrow\; x(ax+b)=0 \;\Rightarrow\; x=0 \quad \text{o} \quad x=-\frac{b}{a}.
\]

Sempre hi ha dues solucions, i una d'elles és $x=0$.
\end{definicio}

\begin{exemple}{\textbf{Exemple:} Equació incompleta sense terme independent}
Resolem $3x^2-12x=0$. Traiem factor comú $3x$:

\[
3x^2-12x=0 \;\Rightarrow\; 3x(x-4)=0.
\]

\begin{itemize}
\item $3x=0 \;\Rightarrow\; x=0$
\item $x-4=0 \;\Rightarrow\; x=4$
\end{itemize}

Les solucions són $x_1=0$ i $x_2=4$.
\end{exemple}

\begin{nota}{Compte: no dividim per x}
Seria un error passar de $3x^2=12x$ a $3x=12$ dividint per $x$: obtindríem només $x=4$ i perdríem la solució $x=0$. Per això sempre traiem factor comú.
\end{nota}

\subsection{Equacions de grau superior}

Per a graus més grans que 2, sovint cal reduir l'equació a una que ja sabem resoldre.

\begin{exemple}{\textbf{Exemple:} Equacions biquadrades}
Una \textbf{equació biquadrada} és una equació de grau 4 de la forma

\[
ax^4+bx^2+c=0.
\]

Per resoldre-la, fem el canvi de variable $t=x^2$. Resolem $x^4-10x^2+9=0$:

\begin{enumerate}
\item Fem el canvi $t=x^2$. L'equació es transforma en

\[
t^2-10t+9=0.
\]
\item Resolem l'equació de segon grau per a $t$:

\[
t=\frac{10\pm\sqrt{100-36}}{2}=\frac{10\pm8}{2}.
\]

Obtenim $t_1=9$ i $t_2=1$.
\item Desfem el canvi de variable per a cada $t$ (recordant que $x=\pm\sqrt{t}$):

\begin{itemize}
\item $t_1=9 \;\Rightarrow\; x^2=9 \;\Rightarrow\; x=\pm\sqrt{9} \;\Rightarrow\; x=\pm 3$
\item $t_2=1 \;\Rightarrow\; x^2=1 \;\Rightarrow\; x=\pm\sqrt{1} \;\Rightarrow\; x=\pm 1$
\end{itemize}
\end{enumerate}

Per tant, les solucions de l'equació són $x_1=3$, $x_2=-3$, $x_3=1$ i $x_4=-1$.
\end{exemple}

\begin{exemple}{\textbf{Exemple:} Equacions de grau superior a 2}
Per a una equació $P(x)=0$ amb $\text{grau}(P)\ge 3$, el mètode habitual és factoritzar el polinomi, ja que les solucions de l'equació són les arrels del polinomi.

Resolem $x^3-3x^2-9x+27=0$:

\begin{enumerate}
\item Busquem arrels enteres (divisors de $27$) i apliquem Ruffini. Provem $x=3$:

\begin{center}
\renewcommand{\arraystretch}{1.25}
\begin{tabular}{r|rrrr}
 & $1$ & $-3$ & $-9$ & $27$ \\
$3$ &  & $3$ & $0$ & $-27$ \\ \hline
 & $1$ & $0$ & $-9$ & $\boxed{0}$ \\
\end{tabular}
\end{center}
\item El polinomi es factoritza com $(x-3)(x^2-9)=0$.
\item Resolem l'equació de segon grau $x^2-9=0$, que dona $x=\pm\sqrt{9}$, és a dir, $x=3$ i $x=-3$.
\item Les solucions són totes les arrels trobades:

\[
x=-3 \ \text{(simple)} \qquad \text{i} \qquad x=3 \ \text{(doble)}.
\]
\end{enumerate}
\end{exemple}

\subsection{Taula resum}

\begin{center}
\begin{tabularx}{\linewidth}{lYY}
\toprule
\textbf{Tipus} & \textbf{Forma} & \textbf{Solucions} \\
\midrule
Primer grau & $ax+b=0,\ a\neq 0$ & $x=-\dfrac{b}{a}$ \\
Segon grau & $ax^2+bx+c=0,\ a\neq 0$ & $x=\dfrac{-b\pm\sqrt{b^2-4ac}}{2a}$ \\
Segon grau incompleta ($b=0$) & $ax^2+c=0$ & $x=\pm\sqrt{-\dfrac{c}{a}}$, si $-\dfrac{c}{a}\geq 0$ \\
Segon grau incompleta ($c=0$) & $ax^2+bx=0$ & $x(ax+b)=0 \Rightarrow x=0,\ x=-\dfrac{b}{a}$ \\
Biquadrada & $ax^4+bx^2+c=0$ & Canvi $t=x^2$, i després $x=\pm\sqrt{t}$ \\
Grau $\geq 3$ & $P(x)=0$ & Factoritzar $P(x)$ (per exemple, amb Ruffini) \\
\bottomrule
\end{tabularx}
\end{center}

\section{Equacions racionals i irracionals}

\subsection{Equacions racionals}

Una equació racional és aquella que involucra fraccions algebraiques.

\begin{definicio}{Definició: equació racional}
Direm que una equació és \textbf{racional} si es pot escriure, amb les transformacions que calguin, com

\[
F(x)=0,
\]

on $F(x)$ és una fracció algebraica. Una manera alternativa d'escriure-la és

\[
\frac{P(x)}{Q(x)} = 0,
\]

on $P(x)$ i $Q(x)$ són polinomis.
\end{definicio}

Les solucions d'aquestes equacions són els valors reals $a$ que anul·len el numerador ($P(a)=0$) però que, alhora, \textbf{no} anul·len el denominador ($Q(a) \neq 0$).

\begin{nota}{Comprovar el domini}
El pas més important en una equació racional és comprovar les solucions.

Qualsevol valor $a$ que faci $Q(a)=0$ (que anul·li el denominador) \textbf{no pot ser solució}, encara que també anul·li el numerador.
\end{nota}

\begin{exemple}{\textbf{Exemple:} Resolució d'una equació racional (Cas 1)}
Resolem l'equació:

\[
\frac{x^2-1}{x-1} = 1.
\]

\begin{enumerate}
\item Transformem l'equació a la forma $\dfrac{P(x)}{Q(x)}=0$:

\[
\frac{x^2-1}{x-1} - 1 = 0 \;\Rightarrow\; \frac{(x^2-1)-(x-1)}{x-1}=0,
\]

\[
\frac{x^2-x}{x-1} = 0.
\]
\item Resolem l'equació del numerador, $P(x)=0$:

\[
x^2-x=0 \;\Rightarrow\; x(x-1)=0.
\]

Les solucions candidates són $x=0$ i $x=1$.
\item Comprovem les solucions al denominador $Q(x)=x-1$:

\begin{itemize}
\item $x=0 \Rightarrow Q(0) = 0-1 = -1 \neq 0$. Per tant, $x=0$ \textbf{és solució}.
\item $x=1 \Rightarrow Q(1) = 1-1 = 0$. Com que anul·la el denominador, $x=1$ \textbf{no és solució}.
\end{itemize}
\end{enumerate}

L'única solució de l'equació és $x=0$.
\end{exemple}

\begin{exemple}{\textbf{Exemple:} Resolució d'una equació racional (Cas 2)}
Resolem l'equació:

\[
\frac{x+1}{x-1}-\frac{4}{x^2-1}=0.
\]

\begin{enumerate}
\item Operem per obtenir una única fracció (recordem que $x^2-1 = (x-1)(x+1)$):

\[
\frac{(x+1)(x+1)}{(x-1)(x+1)} - \frac{4}{x^2-1} = 0,
\]

\[
\frac{(x+1)^2-4}{x^2-1} = 0 \;\Rightarrow\; \frac{x^2+2x+1-4}{x^2-1} = 0,
\]

\[
\frac{x^2+2x-3}{x^2-1} = 0.
\]
\item Resolem $P(x)=0$: $x^2+2x-3=0$. Obtenim les solucions $x_1=1$ i $x_2=-3$.
\item Comprovem les solucions al denominador $Q(x)=x^2-1$:

\begin{itemize}
\item $x_1=1 \Rightarrow Q(1) = 1^2-1 = 0$. Com que anul·la el denominador, \textbf{no és solució}.
\item $x_2=-3 \Rightarrow Q(-3) = (-3)^2-1 = 8 \neq 0$. Per tant, $x=-3$ \textbf{és solució}.
\end{itemize}
\end{enumerate}

L'única solució de l'equació és $x=-3$.
\end{exemple}

\subsection{Equacions irracionals}

\begin{definicio}{Definició: equació irracional}
Direm que una equació és \textbf{irracional} si conté la incògnita sota, almenys, un signe radical. És a dir, si la incògnita està dins d'alguna arrel.
\end{definicio}

El mètode general per resoldre aquestes equacions és aïllar l'arrel (o una de les arrels) i elevar els dos membres de la igualtat al quadrat (o a l'índex de l'arrel) per eliminar el radical.

\begin{nota}{Solucions falses}
Elevar al quadrat els dos membres d'una igualtat \textbf{no és una operació equivalent}:

\[
A=B \;\Rightarrow\; A^2=B^2, \qquad \text{però} \qquad A^2=B^2 \;\Rightarrow\; A=\pm B.
\]

Aquest procés pot introduir solucions que no són vàlides (anomenades \emph{solucions falses}). Per això és \textbf{obligatori comprovar} totes les solucions candidates a l'equació \emph{original}.
\end{nota}

\begin{exemple}{\textbf{Exemple:} Equació irracional amb solució vàlida}
Resolem l'equació:

\[
x + \sqrt{x^2+5} = 1.
\]

\begin{enumerate}
\item Aïllem l'arrel a un costat de la igualtat:

\[
\sqrt{x^2+5} = 1-x.
\]
\item Elevem els dos membres al quadrat:

\[
\left( \sqrt{x^2+5} \right)^2 = (1-x)^2,
\]

\[
x^2+5 = 1 - 2x + x^2.
\]
\item Resolem l'equació resultant (que ja no és irracional):

\[
5 = 1 - 2x \;\Rightarrow\; 2x = 1 - 5 \;\Rightarrow\; 2x = -4.
\]

La solució candidata és $x = -2$.
\item \textbf{Comprovació (obligatòria).} Substituïm $x=-2$ a l'equació \emph{original}:

\[
-2 + \sqrt{(-2)^2 + 5} = -2 + \sqrt{4 + 5} = -2 + \sqrt{9} = -2 + 3 = 1.
\]

Com que $1=1$, la solució $x=-2$ \textbf{és vàlida}.
\end{enumerate}
\end{exemple}

\begin{exemple}{\textbf{Exemple:} Equació irracional amb solució falsa}
Resolem una petita modificació de l'equació anterior:

\[
x - \sqrt{x^2+5} = 1.
\]

\begin{enumerate}
\item Aïllem l'arrel:

\[
-\sqrt{x^2+5} = 1-x.
\]
\item Elevem els dos membres al quadrat:

\[
\left( -\sqrt{x^2+5} \right)^2 = (1-x)^2,
\]

\[
x^2+5 = 1 - 2x + x^2.
\]
\item Resolem l'equació resultant (fixa't que és la mateixa que a l'exemple anterior):

\[
5 = 1 - 2x \;\Rightarrow\; 2x = -4.
\]

La solució candidata és $x = -2$.
\item \textbf{Comprovació (obligatòria).} Substituïm $x=-2$ a l'equació \emph{original}:

\[
-2 - \sqrt{(-2)^2 + 5} = -2 - \sqrt{4 + 5} = -2 - \sqrt{9} = -2 - 3 = -5.
\]

Com que $-5 \neq 1$, la solució $x=-2$ \textbf{no és vàlida} (és una solució falsa). L'equació no té cap solució real.
\end{enumerate}
\end{exemple}

\subsection{Taula resum}

\begin{center}
\begin{tabularx}{\linewidth}{lYY}
\toprule
\textbf{Tipus} & \textbf{Forma} & \textbf{Mètode} \\
\midrule
Racional & $\dfrac{P(x)}{Q(x)}=0$ & Resoldre $P(x)=0$ i descartar les solucions amb $Q(x)=0$ \\
Irracional & Incògnita sota un radical & Aïllar l'arrel, elevar-la a l'índex, i comprovar-ho sempre a l'equació original \\
\bottomrule
\end{tabularx}
\end{center}

\section{Equacions exponencials i logarítmiques}

\subsection{Equacions exponencials}

\begin{definicio}{Definició: equació exponencial}
Una equació és \textbf{exponencial} si apareix alguna potència amb la incògnita a l'exponent.
\end{definicio}

Per resoldre una equació exponencial farem servir tres tècniques bàsiques, sovint combinades:

\begin{enumerate}
\item \textbf{Igualar les bases}, quan podem escriure els dos membres com a potències de la mateixa base.
\item \textbf{Aplicar la definició de logaritme}, quan no és possible igualar les bases.
\item \textbf{Fer un canvi de variable}, quan podem escriure tota l'equació en funció d'una potència $a^x$.
\end{enumerate}

En tots els casos, cal tenir presents les propietats de les potències (com $a^{n+m} = a^n \cdot a^m$ o $(a^n)^m = a^{n\cdot m}$) per transformar les expressions que apareixen a l'equació.

\subsubsection{Igualació de bases}

\begin{propietat}{Propietat: igualació de potències}
Si $a>0$ i $a\neq 1$, aleshores

\[
a^n = a^m \;\Leftrightarrow\; n=m.
\]
\end{propietat}

\begin{exemple}{\textbf{Exemple:} Escrivim les bases com a potències de $2$}
Resolem l'equació $4^{x+1} = 8^x$.

\begin{enumerate}
\item Escrivim $4=2^2$ i $8=2^3$:

\[
\left(2^2\right)^{x+1} = \left(2^3\right)^x.
\]
\item Apliquem la propietat $(a^n)^m=a^{n\cdot m}$:

\[
2^{2x+2} = 2^{3x}.
\]
\item Apliquem la propietat d'igualació de potències:

\[
2x+2 = 3x \;\Rightarrow\; x=2.
\]
\end{enumerate}

Comprovació: $4^{2+1}=4^3=64$ i $8^2=64$.
\end{exemple}

\subsubsection{Aplicació de la definició de logaritme}

En algunes ocasions no és possible igualar les bases. En aquests casos, apliquem la definició de logaritme per aïllar l'exponent.

\begin{propietat}{Propietat: de la potència al logaritme}
Si $a>0$, $a\neq 1$ i $b>0$, aleshores

\[
a^{E(x)} = b \;\Leftrightarrow\; E(x) = \log_a b.
\]

Si $b\leq 0$, l'equació \textbf{no té solució}: una potència de base positiva sempre és positiva. Per exemple, $2^x=-4$ no té solució.
\end{propietat}

\begin{exemple}{\textbf{Exemple:} Resolució aplicant logaritmes}
Resolem l'equació $2^{x+2} = 5$.

\begin{enumerate}
\item Com que $5$ no és una potència clara de $2$, apliquem la definició de logaritme ($\log_a B = C \Leftrightarrow a^C = B$):

\[
x+2 = \log_2 5.
\]
\item Aïllem $x$:

\[
x = -2 + \log_2 5.
\]
\item \emph{(Opcional)} Càlcul i comprovació numèrica. Amb la calculadora, $\log_2 5 = \dfrac{\log 5}{\log 2} \simeq 2{,}3219$:

\[
x \simeq -2 + 2{,}3219 \simeq 0{,}3219.
\]

Si substituïm a l'equació original:

\[
2^{0{,}3219+2} = 2^{2{,}3219} \simeq 4{,}9999\ldots \approx 5.
\]
\end{enumerate}

La solució és $x = -2 + \log_2 5$.
\end{exemple}

\subsubsection{Canvi de variable}

\begin{propietat}{Tècnica: canvi de variable}
Si, aplicant les propietats de les potències, podem escriure tots els termes de l'equació en funció de $a^x$ i de les seves potències ($a^x$, $(a^x)^2$, $(a^x)^3$…), fem el \textbf{canvi de variable} $t=a^x$. Obtenim una equació polinòmica en $t$, que ja sabem resoldre.

Com que $t=a^x$ sempre és positiu, \textbf{descartem les solucions $t\leq 0$\textbf{. Després }desfem el canvi}, igualant bases o aplicant logaritmes.
\end{propietat}

\begin{exemple}{\textbf{Exemple:} Canvi de variable que porta a una equació de primer grau}
Resolem l'equació $3^{x+1} + 3^{x+2} = 36$.

\begin{enumerate}
\item Descomposem les potències usant les propietats:

\[
3^x \cdot 3 + 3^x \cdot 9 = 36.
\]
\item Fem el canvi $t=3^x$:

\[
3t + 9t = 36.
\]
\item Resolem l'equació en $t$:

\[
12t = 36 \;\Rightarrow\; t = 3.
\]
\item Desfem el canvi i apliquem la propietat d'igualació (ja que $3 = 3^1$):

\[
3^x = 3^1 \;\Rightarrow\; x=1.
\]
\end{enumerate}

Fixa't que treure factor comú $3^x$, és a dir, escriure $3^x(3+9)=36$, és fer el mateix sense escriure la $t$.
\end{exemple}

\begin{exemple}{\textbf{Exemple:} Canvi de variable que porta a una equació de segon grau}
Resolem l'equació $4^x - 5\cdot 2^x + 6 = 0$.

\begin{enumerate}
\item Escrivim $4^x$ en funció de $2^x$:

\[
4^x = \left(2^2\right)^x = \left(2^x\right)^2.
\]
\item Fem el canvi $t=2^x$:

\[
t^2 - 5t + 6 = 0.
\]
\item Resolem l'equació de segon grau:

\[
\begin{aligned}
t&=\frac{5\pm\sqrt{25-24}}{2}=\frac{5\pm 1}{2}\\
&\Rightarrow\; t_1=3,\quad t_2=2
\end{aligned}
\]

Tots dos valors són positius, així que no en descartem cap.
\item Desfem el canvi:

\begin{itemize}
\item $2^x = 2 \;\Rightarrow\; x=1$ (igualació de potències).
\item $2^x = 3 \;\Rightarrow\; x=\log_2 3 \simeq 1{,}585$ (definició de logaritme).
\end{itemize}
\end{enumerate}

Comprovació: per a $x=1$, $4-10+6=0$. Per a $x=\log_2 3$, tenim $2^x=3$ i $4^x=\left(2^x\right)^2=9$, i per tant $9-15+6=0$.

Les solucions són $x=1$ i $x=\log_2 3$.
\end{exemple}

\begin{nota}{Compte amb el signe de $t$}
Si resolem $9^x - 2\cdot 3^x - 3 = 0$ amb el canvi $t=3^x$, obtenim $t^2-2t-3=0$, amb solucions $t=3$ i $t=-1$. Descartem $t=-1$, perquè $3^x$ no pot ser negatiu. Només queda $3^x=3$, és a dir, $x=1$.
\end{nota}

\begin{nota}{Quina tècnica faig servir?}
\begin{enumerate}
\item Puc escriure els dos membres com a potències de la mateixa base? \textbf{Igualo els exponents.}
\item Puc escriure tota l'equació en funció de $a^x$ i les seves potències? \textbf{Faig el canvi $t=a^x$.}
\item Em queda $a^{E(x)}=b$ i $b$ no és una potència de $a$? \textbf{Aplico la definició de logaritme.}
\end{enumerate}
\end{nota}

\subsection{Equacions logarítmiques}

\begin{definicio}{Definició: equació logarítmica}
De forma semblant a les equacions exponencials, direm que una equació és \textbf{logarítmica} quan la incògnita apareix dins d'algun logaritme, ja sigui a la base o a l'argument.
\end{definicio}

\begin{nota}{Comprovar les solucions}
L'argument d'un logaritme \textbf{sempre ha de ser positiu}.

És \textbf{obligatori comprovar} totes les solucions obtingudes sobre l'equació \emph{inicial}, i descartar aquelles que facin que un argument sigui zero o negatiu.
\end{nota}

Com abans, caldrà tenir en compte les propietats dels logaritmes per transformar l'equació en una expressió del tipus $\log_b M = \log_b N$. En aquest cas, podrem afirmar que $M=N$.

\begin{exemple}{\textbf{Exemple:} Agrupant termes}
Resolem $\log x - \log 4 = \log(x-2)$.

\begin{enumerate}
\item Apliquem la propietat de la resta de logaritmes:

\[
\log \frac{x}{4} = \log(x-2).
\]
\item Tenim una igualtat de logaritmes de la mateixa base, així que igualem els arguments:

\[
\frac{x}{4} = x-2.
\]
\item Resolem l'equació de primer grau resultant:

\[
x = 4(x-2) \;\Rightarrow\; x = 4x - 8,
\]

\[
3x = 8 \;\Rightarrow\; x = \frac{8}{3}.
\]
\item \textbf{Comprovació.} Verifiquem que tots els arguments de l'equació original són positius:

\begin{itemize}
\item $\log x \Rightarrow \log \tfrac{8}{3}$. Vàlid, $\tfrac{8}{3} > 0$.
\item $\log(x-2) \Rightarrow \log\left(\tfrac{8}{3} - 2\right) = \log \tfrac{2}{3}$. Vàlid, $\tfrac{2}{3} > 0$.
\end{itemize}
\end{enumerate}

Com que tots els arguments són positius, la solució $x=\tfrac{8}{3}$ \textbf{és vàlida}.
\end{exemple}

\begin{exemple}{\textbf{Exemple:} Llei de Hicks-Hyman}
La llei de Hicks-Hyman descriu el temps de reacció ($T$) en funció del nombre d'opcions ($n$):

\[
T=a+b \cdot \log_2 n,
\]

on $a$ i $b$ són constants.

Suposem $a=0{,}6\,\text{s}$ i $b=0{,}1\,\text{s}$. Volem saber quantes opcions $n$ calen per a un temps de reacció de $T=0{,}8\,\text{s}$.

\begin{enumerate}
\item Substituïm els valors coneguts a l'equació:

\[
0{,}8 = 0{,}6 + 0{,}1 \cdot \log_2 n.
\]
\item Aïllem el logaritme:

\[
0{,}8 - 0{,}6 = 0{,}1 \cdot \log_2 n \;\Rightarrow\; 0{,}2 = 0{,}1 \cdot \log_2 n,
\]

\[
\frac{0{,}2}{0{,}1} = \log_2 n \;\Rightarrow\; 2 = \log_2 n.
\]
\item Apliquem la definició de logaritme ($\log_a B = C \Leftrightarrow a^C = B$):

\[
n = 2^2 = 4.
\]
\end{enumerate}

Per tant, per tenir un temps de reacció mitjà de $0{,}8$ segons calen $4$ opcions.
\end{exemple}

\subsection{Taula resum}

\begin{center}
\begin{tabular}{lll}
\toprule
\textbf{Tipus} & \textbf{Propietat clau} & \textbf{Fórmula} \\
\midrule
Exponencial (bases iguals) & Igualació de potències & $a^n=a^m \Leftrightarrow n=m$ \\
Exponencial (bases diferents) & Definició de logaritme & $a^{x}=b \Leftrightarrow x=\log_a b$ \\
Exponencial (en funció de $a^x$) & Canvi de variable & $t=a^x$, amb $t>0$ \\
Logarítmica & Igualació d'arguments & $\log_a M=\log_a N \Leftrightarrow M=N$ \\
\bottomrule
\end{tabular}
\end{center}

\section{Sistemes d'equacions}

\subsection{Equacions amb més d'una incògnita}

Fins ara hem resolt equacions amb una sola incògnita, $x$. Però una equació pot tenir més d'una incògnita, com ara

\[
x+y=5, \qquad x^2+y^2=25, \qquad xy=6.
\]

\begin{definicio}{Definició: solució d'una equació amb diverses incògnites}
Una \textbf{solució} d'una equació amb diverses incògnites és un conjunt de valors, \textbf{un per a cada incògnita}, que fan certa la igualtat.

Quan les incògnites són $x$ i $y$, escrivim una solució com el parell $(x,y)$, amb el valor de $x$ primer i el de $y$ després.
\end{definicio}

\begin{exemple}{\textbf{Exemple:} l'equació $x+y=5$}
El parell $(2,3)$ és solució, perquè $2+3=5$. També ho són $(0,5)$, $(-1,6)$ o $\left(\dfrac12,\dfrac92\right)$.

De fet, triem el valor de $x$ que triem, n'hi ha prou de prendre $y=5-x$ per obtenir una solució. Per tant, l'equació té \textbf{infinites solucions}.
\end{exemple}

En general, una equació amb dues o més incògnites té infinites solucions. Hi ha excepcions: per exemple, $x^2+y^2=0$ només té la solució $(0,0)$, i $x^2+y^2=-1$ no en té cap, perquè una suma de quadrats no pot ser negativa.

\subsection{Sistemes d'equacions}

\begin{definicio}{Definició: sistema d'equacions}
Un \textbf{sistema d'equacions} és un conjunt d'equacions amb les mateixes incògnites, de les quals busquem les solucions comunes.

Una \textbf{solució del sistema} és un conjunt de valors que és solució de \textbf{totes} les equacions alhora. \textbf{Resoldre} el sistema vol dir trobar totes les seves solucions o bé veure que no en té cap.
\end{definicio}

Per exemple, el parell $(4,3)$ és solució del sistema

\[
\begin{cases}
x-y=1\\
x^2+y^2=25
\end{cases}
\]

perquè compleix les dues equacions. Substituïm $x=4$ i $y=3$ a cada una:

\[
\begin{aligned}
\text{1a equació: }&4-3=1\\
\text{2a equació: }&4^2+3^2=16+9=25
\end{aligned}
\]

En canvi, el parell $(3,4)$ compleix la segona equació, però no la primera:

\[
\begin{aligned}
\text{1a equació: }&3-4=-1\neq 1\\
\text{2a equació: }&3^2+4^2=9+16=25
\end{aligned}
\]

Per tant, $(3,4)$ no és solució del sistema.

Ja coneixeu els \textbf{sistemes lineals}, en què totes les equacions són de primer grau, de dues equacions amb dues incògnites o de tres equacions amb tres incògnites. Ara treballarem amb \textbf{sistemes no lineals}: sistemes en què alguna de les equacions no és de primer grau.

\begin{exemple}{\textbf{Exemple:} sistemes no lineals}
En aquests sistemes hi apareixen quadrats, productes d'incògnites, fraccions, arrels, potències o logaritmes:

\[
\begin{cases}
x+y=7\\
xy=12
\end{cases}
\qquad
\begin{cases}
y=x^2-4x+3\\
y=x-1
\end{cases}
\]

\[
\begin{cases}
x-y=1\\
\sqrt{x+y}+y=7
\end{cases}
\qquad
\begin{cases}
\log x+\log y=2\\
x-y=15
\end{cases}
\]
\end{exemple}

\begin{nota}{Tantes equacions com incògnites}
Treballarem amb sistemes que tenen \textbf{tantes equacions com incògnites}, sobretot de dues equacions amb dues incògnites.

Una sola equació amb dues incògnites sol tenir infinites solucions. Quan n'hi afegim una segona, les solucions han de complir totes dues equacions alhora, i això en redueix molt el nombre: els sistemes que veurem tenen \textbf{un nombre finit de solucions} (una, dues, quatre...) o \textbf{cap}. Només hi ha infinites solucions en casos especials, per exemple quan una equació és equivalent a l'altra, com a $x+y=2$ i $2x+2y=4$ (la segona és la primera multiplicada per $2$). En aquests casos és com si només tinguéssim una equació, i per això torna a haver-hi infinites solucions.
\end{nota}

\subsection{Mètodes de resolució}

Els sistemes no lineals es resolen amb els mateixos mètodes que els lineals. Nosaltres farem servir la \textbf{substitució} i la \textbf{igualació}. La idea és sempre la mateixa: obtenir una \textbf{equació amb una sola incògnita}, que ja sabem resoldre.

\begin{propietat}{Mètode de substitució}
\begin{enumerate}
\item Aïllem una incògnita en una de les equacions (la que sigui més senzilla).
\item Substituïm l'expressió obtinguda a l'altra equació, que queda amb una sola incògnita.
\item Resolem aquesta equació.
\item Per a cada valor trobat, calculem l'altra incògnita amb l'expressió del pas 1.
\end{enumerate}
\end{propietat}

\begin{propietat}{Mètode d'igualació}
\begin{enumerate}
\item Aïllem la mateixa incògnita a les dues equacions.
\item Igualem les dues expressions: obtenim una equació amb una sola incògnita.
\item Resolem aquesta equació.
\item Per a cada valor trobat, calculem l'altra incògnita amb qualsevol de les dues expressions.
\end{enumerate}
\end{propietat}

\begin{nota}{Parells de solucions i comprovació}
\begin{itemize}
\item Cada valor d'una incògnita dona el seu valor de l'altra: les solucions s'escriuen per \textbf{parells} $(x,y)$, i no es poden barrejar els valors d'un parell amb els d'un altre.
\item Si en el procés elevem al quadrat, multipliquem per una expressió amb incògnites o fem servir logaritmes, poden aparèixer solucions falses. Cal \textbf{comprovar} les solucions al sistema original, igual que fèiem amb les equacions.
\end{itemize}
\end{nota}

\subsection{Exemples}

\begin{exemple}{\textbf{Exemple 1:} equació de segon grau (substitució)}
Resolem el sistema

\[
\begin{cases}
x-y=1\\
x^2+y^2=25
\end{cases}
\]

\begin{enumerate}
\item Aïllem $y$ a la primera equació: $y=x-1$.
\item Substituïm a la segona equació i resolem:

\[
\begin{aligned}
x^2+(x-1)^2&=25\\
x^2+x^2-2x+1&=25\\
2x^2-2x-24&=0\\
x^2-x-12&=0
\end{aligned}
\]

\[
\begin{aligned}x&=\frac{1\pm\sqrt{1+48}}{2}=\frac{1\pm 7}{2}\\ &\Rightarrow\; x_1=4,\quad x_2=-3\end{aligned}
\]
\item Calculem $y=x-1$ per a cada valor:

\begin{itemize}
\item $x_1=4 \;\Rightarrow\; y_1=3$
\item $x_2=-3 \;\Rightarrow\; y_2=-4$
\end{itemize}
\end{enumerate}

El sistema té dues solucions: $(4,3)$ i $(-3,-4)$.
\end{exemple}

\begin{exemple}{\textbf{Exemple 2:} una paràbola i una recta (igualació)}
Resolem el sistema

\[
\begin{cases}
y=x^2-4x+3\\
y=x-1
\end{cases}
\]

\begin{enumerate}
\item Les dues equacions ja tenen $y$ aïllada. Igualem les dues expressions:

\[
\begin{aligned}
x^2-4x+3&=x-1\\
x^2-5x+4&=0
\end{aligned}
\]

\[
\begin{aligned}x&=\frac{5\pm\sqrt{25-16}}{2}=\frac{5\pm 3}{2}\\ &\Rightarrow\; x_1=4,\quad x_2=1\end{aligned}
\]
\item Calculem $y=x-1$ per a cada valor:

\begin{itemize}
\item $x_1=4 \;\Rightarrow\; y_1=3$
\item $x_2=1 \;\Rightarrow\; y_2=0$
\end{itemize}
\end{enumerate}

El sistema té dues solucions: $(4,3)$ i $(1,0)$.

Gràficament, la primera equació és una paràbola i la segona una recta, i les solucions del sistema són els seus \textbf{punts de tall}, $(1,0)$ i $(4,3)$. Ho pots veure al Laboratori d'equacions escrivint $x^2-4x+3$ i $x-1$ als dos membres: les abscisses dels punts de tall són les solucions de $x^2-4x+3=x-1$.
\end{exemple}

\begin{exemple}{\textbf{Exemple 3:} un producte d'incògnites i una biquadrada (substitució)}
Resolem el sistema

\[
\begin{cases}
x^2+y^2=13\\
xy=6
\end{cases}
\]

\begin{enumerate}
\item Aïllem $y$ a la segona equació: $y=\dfrac{6}{x}$. Podem dividir per $x$ perquè $x\neq 0$: si fos $x=0$, la segona equació diria $0=6$.
\item Substituïm a la primera equació i multipliquem per $x^2$:

\[
\begin{aligned}
x^2+\frac{36}{x^2}&=13\\
x^4+36&=13x^2\\
x^4-13x^2+36&=0
\end{aligned}
\]

És una equació biquadrada. Amb el canvi $t=x^2$:

\[
\begin{aligned}t^2-13t+36=0 \;&\Rightarrow\; t=\frac{13\pm 5}{2}\\ &\Rightarrow\; t_1=9,\quad t_2=4\end{aligned}
\]

i, desfent el canvi, $x=\pm 3$ i $x=\pm 2$.
\item Calculem $y=\dfrac{6}{x}$ per a cada valor:

\begin{itemize}
\item $x=3 \Rightarrow y=2$, $\quad x=-3 \Rightarrow y=-2$
\item $x=2 \Rightarrow y=3$, $\quad x=-2 \Rightarrow y=-3$
\end{itemize}
\end{enumerate}

El sistema té quatre solucions: $(3,2)$, $(-3,-2)$, $(2,3)$ i $(-2,-3)$.
\end{exemple}

\begin{exemple}{\textbf{Exemple 4:} una equació irracional (substitució)}
Resolem el sistema

\[
\begin{cases}
x-y=1\\
\sqrt{x+y}+y=7
\end{cases}
\]

\begin{enumerate}
\item Aïllem $x$ a la primera equació: $x=y+1$.
\item Substituïm a la segona equació. Aïllem l'arrel i elevem al quadrat:

\[
\begin{aligned}
\sqrt{2y+1}+y&=7\\
\sqrt{2y+1}&=7-y\\
2y+1&=49-14y+y^2\\
y^2-16y+48&=0
\end{aligned}
\]

\[
\begin{aligned}y&=\frac{16\pm\sqrt{256-192}}{2}=\frac{16\pm 8}{2}\\ &\Rightarrow\; y_1=12,\quad y_2=4\end{aligned}
\]
\item Calculem $x=y+1$: obtenim els parells $(13,12)$ i $(5,4)$.
\item \textbf{Comprovació.} Hem elevat al quadrat, així que comprovem les dues candidates a la segona equació, $\sqrt{x+y}+y=7$ (la primera es compleix sempre, perquè hem calculat $x=y+1$):

\[
\begin{aligned}
(13,12)\text{: }&\sqrt{13+12}+12=5+12=17\neq 7\\
(5,4)\text{: }&\sqrt{5+4}+4=3+4=7
\end{aligned}
\]

El parell $(13,12)$ és una solució falsa, que ha aparegut en elevar al quadrat. El parell $(5,4)$ sí que és solució.
\end{enumerate}

El sistema té una única solució: $(5,4)$.
\end{exemple}

\begin{exemple}{\textbf{Exemple 5:} equacions exponencials}
Resolem el sistema

\[
\begin{cases}
2^x\cdot 2^y=32\\
3^x=3\cdot 3^y
\end{cases}
\]

\begin{enumerate}
\item Transformem cada equació amb les propietats de les potències i la igualació de potències:

\[
\begin{aligned}
2^x\cdot 2^y=2^{x+y}=2^5 &\;\Rightarrow\; x+y=5\\
3^x=3^{1+y} &\;\Rightarrow\; x=1+y
\end{aligned}
\]
\item Ens queda un sistema lineal. Substituïm $x=1+y$ a la primera equació:

\[
1+y+y=5 \;\Rightarrow\; 2y=4 \;\Rightarrow\; y=2
\]

i, per tant, $x=1+2=3$.
\end{enumerate}

El sistema té una única solució: $(3,2)$.
\end{exemple}

\begin{exemple}{\textbf{Exemple 6:} equacions logarítmiques (substitució)}
Resolem el sistema

\[
\begin{cases}
\log x+\log y=2\\
x-y=15
\end{cases}
\]

\begin{enumerate}
\item A la primera equació apliquem la propietat del producte i la definició de logaritme:

\[
\log(xy)=2 \;\Rightarrow\; xy=10^2=100
\]
\item Aïllem $x$ a la segona equació, $x=y+15$, i substituïm:

\[
\begin{aligned}
(y+15)\,y&=100\\
y^2+15y-100&=0
\end{aligned}
\]

\[
\begin{aligned}y&=\frac{-15\pm\sqrt{225+400}}{2}\\ &=\frac{-15\pm 25}{2}\\ &\Rightarrow\; y_1=5,\quad y_2=-20\end{aligned}
\]
\item Calculem $x=y+15$: obtenim els parells $(20,5)$ i $(-5,-20)$.
\item \textbf{Comprovació.} Els arguments dels logaritmes han de ser positius. Substituïm cada parell a la primera equació, $\log x+\log y=2$:

\[
\begin{aligned}
(20,5)\text{: }&\log 20+\log 5=\log 100=2\\
(-5,-20)\text{: }&\log(-5)\ \text{no existeix}
\end{aligned}
\]

El parell $(20,5)$ és solució (també compleix $20-5=15$). El parell $(-5,-20)$ no ho és.
\end{enumerate}

El sistema té una única solució: $(20,5)$.
\end{exemple}

\subsection{Taula resum}

\begin{center}
\begin{tabularx}{\linewidth}{lYY}
\toprule
\textbf{Mètode} & \textbf{Com es fa} & \textbf{Quan és útil} \\
\midrule
Substitució & S'aïlla una incògnita en una equació i se substitueix a l'altra & Quan una incògnita és fàcil d'aïllar (equacions de primer grau, $xy=k$...) \\
Igualació & S'aïlla la mateixa incògnita a les dues equacions i s'igualen & Quan la mateixa incògnita ja és aïllada, o és fàcil d'aïllar, a les dues equacions \\
En tots dos casos & S'obté una equació amb una incògnita; es calculen els parells $(x,y)$ & Cal comprovar les solucions si s'ha elevat al quadrat, s'ha multiplicat per una expressió amb incògnites o hi ha logaritmes \\
\bottomrule
\end{tabularx}
\end{center}

\section{Inequacions}

\subsection{Què és una inequació}

\begin{definicio}{Definició: inequació}
Una \textbf{inequació} és una desigualtat entre dues expressions algebraiques en què apareix una incògnita. Fem servir els signes $<$, $\leq$, $>$ i $\geq$.

Una \textbf{solució} d'una inequació és un valor de la incògnita que fa certa la desigualtat. \textbf{Resoldre} la inequació vol dir trobar \textbf{totes} les solucions.
\end{definicio}

\begin{exemple}{\textbf{Exemple:} la inequació $2x+1>7$}
Provem alguns valors:

\[
\begin{aligned}
x=5&:\quad 2\cdot 5+1=11>7 &&\text{(és solució)}\\
x=3&:\quad 2\cdot 3+1=7\not>7 &&\text{(no és solució)}\\
x=1&:\quad 2\cdot 1+1=3\not>7 &&\text{(no és solució)}
\end{aligned}
\]

De fet, són solució \textbf{tots} els nombres més grans que $3$. La inequació té infinites solucions, que escrivim amb un interval: $(3,+\infty)$.
\end{exemple}

A diferència de les equacions, que normalment tenen un nombre finit de solucions, les inequacions solen tenir-ne infinites. Per això les escrivim amb intervals o unions d'intervals. Recordem que:

\begin{itemize}
\item Amb $<$ o $>$, l'extrem \textbf{no} hi és inclòs i fem servir un parèntesi: $x<4 \Rightarrow (-\infty,4)$.
\item Amb $\leq$ o $\geq$, l'extrem \textbf{sí} que hi és inclòs i fem servir un claudàtor: $x\leq 4 \Rightarrow (-\infty,4]$.
\item Per als extrems $-\infty$ i $+\infty$ sempre fem servir parèntesis.
\end{itemize}

\subsection{Inequacions de primer grau}

Les inequacions de primer grau es resolen com les equacions de primer grau: traiem parèntesis i denominadors, agrupem els termes amb $x$ en un membre i aïllem la $x$. Les transformacions que podem fer són:

\begin{itemize}
\item \textbf{Sumar o restar} el mateix nombre o expressió als dos membres: la desigualtat \textbf{no canvia}.
\item \textbf{Multiplicar o dividir} els dos membres per un nombre \textbf{positiu}: la desigualtat \textbf{no canvia}.
\item \textbf{Multiplicar o dividir} els dos membres per un nombre \textbf{negatiu}: la desigualtat \textbf{canvia de sentit}.
\end{itemize}

\begin{compte}{Compte: multiplicar o dividir per un nombre negatiu}
Sabem que $2<3$. Si multipliquem els dos nombres per $-1$, obtenim $-2$ i $-3$, i ara és $-2$ el més gran:

\[
2<3 \quad\Rightarrow\quad -2>-3
\]

A la recta real, canviar el signe és com mirar-se en un mirall col·locat al $0$: l'ordre dels nombres s'inverteix. Per això, \textbf{quan multipliquem o dividim per un nombre negatiu, hem de girar la desigualtat}: $<$ passa a ser $>$, i $\leq$ passa a ser $\geq$ (i al revés).
\end{compte}

\begin{exemple}{\textbf{Exemple 1:} $3x-2\leq 10$}
Passem el $-2$ a l'altre membre i dividim per $3$:

\[
\begin{aligned}
3x-2&\leq 10\\
3x&\leq 12\\
x&\leq \frac{12}{3}=4
\end{aligned}
\]

Hem dividit per $3$, que és \textbf{positiu}, i per tant la desigualtat \textbf{no} canvia. La solució és

\[
(-\infty,4].
\]

Comprovem-ho amb un valor de dins de l'interval i un de fora:

\[
\begin{aligned}
x=0&:\quad 3\cdot 0-2=-2\leq 10 &&\text{(és solució)}\\
x=5&:\quad 3\cdot 5-2=13\not\leq 10 &&\text{(no és solució)}
\end{aligned}
\]
\end{exemple}

\begin{exemple}{\textbf{Exemple 2:} $2(x-3)\leq 5x+3$}
Traiem el parèntesi i agrupem:

\[
\begin{aligned}
2x-6&\leq 5x+3\\
2x-5x&\leq 3+6\\
-3x&\leq 9
\end{aligned}
\]

Ara hem de dividir per $-3$, que és \textbf{negatiu}: \textbf{girem la desigualtat}.

\[
x\geq \frac{9}{-3}=-3
\]

La solució és

\[
[-3,+\infty).
\]

Comprovem-ho amb un valor de dins de l'interval i un de fora:

\[
\begin{aligned}
x=0&:\quad 2(0-3)\overset{?}{\leq} 5\cdot 0+3\\
&\Rightarrow\; -6\leq 3 \quad\text{(és solució)}\\[6pt]
x=-4&:\quad 2(-4-3)\overset{?}{\leq} 5(-4)+3\\
&\Rightarrow\; -14\not\leq -17 \quad\text{(no és solució)}
\end{aligned}
\]
\end{exemple}

\begin{propietat}{Per no haver de girar la desigualtat}
Podem agrupar els termes amb $x$ al membre on el coeficient queda positiu. A l'exemple 2:

\[
\begin{aligned}
2x-6&\leq 5x+3\\
-6-3&\leq 5x-2x\\
-9&\leq 3x\\
-3&\leq x
\end{aligned}
\]

Obtenim el mateix resultat, $x\geq -3$, sense haver dividit per cap nombre negatiu.
\end{propietat}

\subsubsection{Una altra manera de mirar-ho}

Tornem a la inequació de l'exemple 1, $3x-2\leq 10$. Si passem el $10$ a l'esquerra, obtenim la inequació equivalent

\[
3x-12\leq 0.
\]

Si dibuixem la recta $y=3x-12$, les solucions són els valors de $x$ on la recta és \textbf{per sota} de l'eix $OX$ o el toca:

\begin{nota}{Com llegim les gràfiques d'aquest apartat}
\begin{itemize}
\item \mostracorba{} Tros de la gràfica que \textbf{compleix} la desigualtat.
\item \mostrainterval{} \textbf{Solució}, marcada sobre l'eix $OX$.
\item \mostraple{} Punt \textbf{ple}: l'extrem \textbf{hi és inclòs} ($\leq$ o $\geq$).
\item \mostrabuit{} Punt \textbf{buit}: l'extrem \textbf{no hi és inclòs} ($<$ o $>$).
\end{itemize}
\end{nota}

\begin{center}
\begin{tikzpicture}
\begin{axis}[grafica, width=0.62\linewidth, height=0.5*0.62\linewidth, xmin=-1.5, xmax=7.5, ymin=-14, ymax=8]
\addplot[g-corba, domain=-1.5:7.5, samples=160] {0.0*x^2 + 3.0*x + -12.0};
\addplot[g-solcorba, domain=-1.5:4, samples=100] {0.0*x^2 + 3.0*x + -12.0};
\draw[g-interval] (axis cs:-1.5,0) -- (axis cs:4,0);
\draw[g-ple] (axis cs:4,0) circle[radius=2.4pt];
\node[below=3pt, font=\scriptsize, fill=white, inner sep=1pt] at (axis cs:4,0) {$4$};
\end{axis}
\end{tikzpicture}
\end{center}

Aquesta idea, mirar on la gràfica és per sobre o per sota de l'eix $OX$, és la que farem servir per resoldre les inequacions de segon grau.

\subsection{Inequacions de segon grau}

Una inequació és de \textbf{segon grau} quan, passant tots els termes a un membre i simplificant, la podem escriure en la \textbf{forma canònica}

\[
ax^2+bx+c<0, \qquad a\neq 0,
\]

o bé la mateixa expressió amb $\leq$, $>$ o $\geq$.

El primer pas sempre és portar la inequació a aquesta forma, deixant un $0$ en un dels membres. Per exemple:

\[
2x^2\geq x+3 \quad\Rightarrow\quad 2x^2-x-3\geq 0.
\]

\begin{definicio}{Idea clau}
La gràfica de $y=ax^2+bx+c$ és una \textbf{paràbola}. Per tant:

\begin{itemize}
\item $ax^2+bx+c>0$ vol dir buscar on la paràbola és \textbf{per sobre} de l'eix $OX$.
\item $ax^2+bx+c<0$ vol dir buscar on la paràbola és \textbf{per sota} de l'eix $OX$.
\item Amb $\geq$ o $\leq$, també hi afegim els punts on la paràbola \textbf{talla} l'eix $OX$.
\end{itemize}
\end{definicio}

Per saber on és per sobre i on és per sota, en tenim prou de conèixer dues coses: la \textbf{forma} de la paràbola i els seus \textbf{punts de tall} amb l'eix $OX$.

\subsubsection{La forma de la paràbola}

\begin{definicio}{Forma de la paràbola segons el signe de $a$}
\begin{itemize}
\item Si $a>0$, la paràbola és \textbf{còncava} ($\cup$): és oberta cap amunt i les dues branques pugen.
\item Si $a<0$, la paràbola és \textbf{convexa} ($\cap$): és oberta cap avall i les dues branques baixen.
\end{itemize}
\end{definicio}

\begin{center}
\begin{minipage}[t]{0.48\linewidth}
\centering
\begin{center}
\begin{tikzpicture}
\begin{axis}[grafica, width=0.92\linewidth, height=0.5*0.92\linewidth, xmin=-3, xmax=5, ymin=-5, ymax=6]
\addplot[g-corba, domain=-3:5, samples=160] {1.0*x^2 + -2.0*x + -3.0};
\fill[g-arrel] (axis cs:-1,0) circle[radius=2.2pt];
\node[below=3pt, font=\scriptsize, fill=white, inner sep=1pt] at (axis cs:-1,0) {$-1$};
\fill[g-arrel] (axis cs:3,0) circle[radius=2.2pt];
\node[below=3pt, font=\scriptsize, fill=white, inner sep=1pt] at (axis cs:3,0) {$3$};
\end{axis}
\end{tikzpicture}
\end{center}

$y=x^2-2x-3$: $a=1>0$, còncava ($\cup$)
\end{minipage}\hfill
\begin{minipage}[t]{0.48\linewidth}
\centering
\begin{center}
\begin{tikzpicture}
\begin{axis}[grafica, width=0.92\linewidth, height=0.5*0.92\linewidth, xmin=-3, xmax=5, ymin=-6, ymax=5]
\addplot[g-corba, domain=-3:5, samples=160] {-1.0*x^2 + 2.0*x + 3.0};
\fill[g-arrel] (axis cs:-1,0) circle[radius=2.2pt];
\node[below=3pt, font=\scriptsize, fill=white, inner sep=1pt] at (axis cs:-1,0) {$-1$};
\fill[g-arrel] (axis cs:3,0) circle[radius=2.2pt];
\node[below=3pt, font=\scriptsize, fill=white, inner sep=1pt] at (axis cs:3,0) {$3$};
\end{axis}
\end{tikzpicture}
\end{center}

$y=-x^2+2x+3$: $a=-1<0$, convexa ($\cap$)
\end{minipage}
\end{center}

Per què depèn només de $a$? Quan $x$ és molt gran (positiu o negatiu), el terme $ax^2$ és molt més gran que els altres i decideix el signe. Com que $x^2$ sempre és positiu, el signe de $ax^2+bx+c$ per a valors grans de $x$ és el signe de $a$.

\begin{nota}{Compte amb els noms}
No tots els llibres fan servir el mateix conveni: n'hi ha que diuen \emph{convexa} a la forma $\cup$ i \emph{còncava} a la forma $\cap$. Per evitar confusions, sempre podem dir \textbf{oberta cap amunt} o \textbf{oberta cap avall}.
\end{nota}

\subsubsection{\texorpdfstring{Els punts de tall amb l'eix $OX$}{Els punts de tall amb l'eix OX}}

Els punts de tall són les solucions de l'equació $ax^2+bx+c=0$. Depenen del signe del discriminant $\Delta=b^2-4ac$:

\begin{itemize}
\item Si $\Delta>0$, la paràbola talla l'eix $OX$ en \textbf{dos punts}.
\item Si $\Delta<0$, la paràbola \textbf{no talla} l'eix $OX$.
\item Si $\Delta=0$, la paràbola toca l'eix $OX$ en \textbf{un sol punt}: hi és \textbf{tangent}.
\end{itemize}

Estudiem cada cas per separat.

\subsection{Cas general: dos punts de tall}

Quan $\Delta>0$, l'equació $ax^2+bx+c=0$ té dues solucions, $x_1<x_2$, i la recta real queda dividida en tres trossos.

\begin{definicio}{Pas a pas}
\begin{enumerate}
\item Portem la inequació a la \textbf{forma canònica}: $ax^2+bx+c$ comparat amb $0$.
\item Mirem el \textbf{signe de $a$}: si $a>0$, la paràbola és $\cup$; si $a<0$, és $\cap$.
\item Resolem l'equació $ax^2+bx+c=0$ per trobar els \textbf{punts de tall} $x_1<x_2$.
\item Fem un \textbf{esbós} de la paràbola: talla l'eix $OX$ a $x_1$ i $x_2$ i té la forma del pas 2.
\item Fem la \textbf{taula de signes}: amb els punts de tall dividim la recta en intervals i hi anotem el signe que llegim a l'esbós ($+$ si la paràbola és per sobre de l'eix, $-$ si és per sota i $0$ als punts de tall). Si no tenim l'esbós, trobem el signe de cada interval calculant el polinomi en un valor qualsevol de dins.
\item \textbf{Llegim la solució} a la taula: per a $>0$, els intervals amb $+$; per a $<0$, els intervals amb $-$.
\item Decidim els \textbf{extrems}: amb $\geq$ o $\leq$, els punts de tall (amb signe $0$) hi són inclosos i fem servir claudàtors; amb $>$ o $<$, no hi són inclosos i fem servir parèntesis.
\end{enumerate}
\end{definicio}

Si $a>0$ ($\cup$), la paràbola és per sota de l'eix \textbf{entre} les arrels i per sobre \textbf{fora} de les arrels. Si $a<0$ ($\cap$), és al revés.

\begin{exemple}{\textbf{Exemple 1:} $x^2-x-6<0$}
Ja és en forma canònica. Com que $a=1>0$, la paràbola és \textbf{còncava} ($\cup$). Busquem els punts de tall:

\[
\begin{aligned}
x^2-x-6&=0\\
x&=\frac{1\pm\sqrt{1+24}}{2}=\frac{1\pm 5}{2}\\
x_1&=-2,\quad x_2=3
\end{aligned}
\]

\begin{center}
\begin{tikzpicture}
\begin{axis}[grafica, width=0.62\linewidth, height=0.5*0.62\linewidth, xmin=-4, xmax=5, ymin=-8, ymax=8]
\addplot[g-corba, domain=-4:5, samples=160] {1.0*x^2 + -1.0*x + -6.0};
\addplot[g-solcorba, domain=-2:3, samples=100] {1.0*x^2 + -1.0*x + -6.0};
\draw[g-interval] (axis cs:-2,0) -- (axis cs:3,0);
\draw[g-buit] (axis cs:-2,0) circle[radius=2.4pt];
\node[below=3pt, font=\scriptsize, fill=white, inner sep=1pt] at (axis cs:-2,0) {$-2$};
\draw[g-buit] (axis cs:3,0) circle[radius=2.4pt];
\node[below=3pt, font=\scriptsize, fill=white, inner sep=1pt] at (axis cs:3,0) {$3$};
\end{axis}
\end{tikzpicture}
\end{center}

Fem la taula de signes. Els signes es llegeixen a la gràfica; també els podem trobar amb un valor de prova a cada interval:

\begin{center}
\begin{tabular}{lccccc}
\toprule
$x$ & $(-\infty,-2)$ & $-2$ & $(-2,3)$ & $3$ & $(3,+\infty)$ \\
\midrule
Prova & $-3$ &  & $0$ &  & $4$ \\
$x^2-x-6$ & $6$ & $0$ & $-6$ & $0$ & $6$ \\
Signe & $+$ & $0$ & $-$ & $0$ & $+$ \\
$<0$? & no & no & \textbf{sí} & no & no \\
\bottomrule
\end{tabular}
\end{center}

Volem que sigui $<0$: ens quedem amb l'interval amb signe $-$, \textbf{entre} les arrels. Com que la desigualtat és estricta, els punts de tall (signe $0$) no hi són inclosos:

\[
(-2,3)
\]
\end{exemple}

\begin{exemple}{\textbf{Exemple 2:} $2x^2\geq x+3$}
Primer la portem a la forma canònica:

\[
2x^2-x-3\geq 0
\]

Com que $a=2>0$, la paràbola és \textbf{còncava} ($\cup$). Busquem els punts de tall:

\[
\begin{aligned}
2x^2-x-3&=0\\
x&=\frac{1\pm\sqrt{1+24}}{4}=\frac{1\pm 5}{4}\\
x_1&=-1,\quad x_2=\frac32
\end{aligned}
\]

\begin{center}
\begin{tikzpicture}
\begin{axis}[grafica, width=0.62\linewidth, height=0.5*0.62\linewidth, xmin=-3, xmax=3.5, ymin=-5, ymax=6]
\addplot[g-corba, domain=-3:3.5, samples=160] {2.0*x^2 + -1.0*x + -3.0};
\addplot[g-solcorba, domain=-3:-1, samples=100] {2.0*x^2 + -1.0*x + -3.0};
\draw[g-interval] (axis cs:-3,0) -- (axis cs:-1,0);
\addplot[g-solcorba, domain=1.5:3.5, samples=100] {2.0*x^2 + -1.0*x + -3.0};
\draw[g-interval] (axis cs:1.5,0) -- (axis cs:3.5,0);
\draw[g-ple] (axis cs:-1,0) circle[radius=2.4pt];
\node[below=3pt, font=\scriptsize, fill=white, inner sep=1pt] at (axis cs:-1,0) {$-1$};
\draw[g-ple] (axis cs:1.5,0) circle[radius=2.4pt];
\node[below=3pt, font=\scriptsize, fill=white, inner sep=1pt] at (axis cs:1.5,0) {$3/2$};
\end{axis}
\end{tikzpicture}
\end{center}

Fem la taula de signes:

\begin{center}
\begin{tabular}{lccccc}
\toprule
$x$ & $(-\infty,-1)$ & $-1$ & $(-1,\,3/2)$ & $3/2$ & $(3/2,+\infty)$ \\
\midrule
Signe & $+$ & $0$ & $-$ & $0$ & $+$ \\
$\geq 0$? & \textbf{sí} & \textbf{sí} & no & \textbf{sí} & \textbf{sí} \\
\bottomrule
\end{tabular}
\end{center}

Volem que sigui $\geq 0$: els intervals amb signe $+$, \textbf{fora} de les arrels, i també els punts de tall (signe $0$), perquè la desigualtat inclou l'igual:

\[
(-\infty,-1]\cup\left[\frac32,+\infty\right)
\]

Comprovem-ho amb un valor de dins i un de fora:

\[
\begin{aligned}
x=2&:\quad 2\cdot 2^2=8\geq 2+3=5 &&\text{(és solució)}\\
x=0&:\quad 2\cdot 0^2=0\not\geq 0+3=3 &&\text{(no és solució)}
\end{aligned}
\]
\end{exemple}

\begin{exemple}{\textbf{Exemple 3:} $-x^2+4x-3>0$}
Ara $a=-1<0$: la paràbola és \textbf{convexa} ($\cap$). Busquem els punts de tall:

\[
\begin{aligned}
-x^2+4x-3&=0\\
x^2-4x+3&=0\\
x&=\frac{4\pm\sqrt{16-12}}{2}=\frac{4\pm 2}{2}\\
x_1&=1,\quad x_2=3
\end{aligned}
\]

\begin{center}
\begin{tikzpicture}
\begin{axis}[grafica, width=0.62\linewidth, height=0.5*0.62\linewidth, xmin=-1, xmax=5, ymin=-5, ymax=3]
\addplot[g-corba, domain=-1:5, samples=160] {-1.0*x^2 + 4.0*x + -3.0};
\addplot[g-solcorba, domain=1:3, samples=100] {-1.0*x^2 + 4.0*x + -3.0};
\draw[g-interval] (axis cs:1,0) -- (axis cs:3,0);
\draw[g-buit] (axis cs:1,0) circle[radius=2.4pt];
\node[below=3pt, font=\scriptsize, fill=white, inner sep=1pt] at (axis cs:1,0) {$1$};
\draw[g-buit] (axis cs:3,0) circle[radius=2.4pt];
\node[below=3pt, font=\scriptsize, fill=white, inner sep=1pt] at (axis cs:3,0) {$3$};
\end{axis}
\end{tikzpicture}
\end{center}

Fem la taula de signes. Com que la paràbola és oberta cap avall, el signe $+$ ara és \textbf{entre} les arrels:

\begin{center}
\begin{tabular}{lccccc}
\toprule
$x$ & $(-\infty,1)$ & $1$ & $(1,3)$ & $3$ & $(3,+\infty)$ \\
\midrule
Signe & $-$ & $0$ & $+$ & $0$ & $-$ \\
$>0$? & no & no & \textbf{sí} & no & no \\
\bottomrule
\end{tabular}
\end{center}

Volem que sigui $>0$: l'interval amb signe $+$, sense els extrems:

\[
(1,3)
\]

Comprovem-ho: per a $x=2$ tenim $-4+8-3=1>0$, que és solució; per a $x=0$ tenim $-3\not>0$, que no ho és.
\end{exemple}

\begin{propietat}{Una altra manera: fer que $a$ sigui positiu}
Si multipliquem per $-1$ (i \textbf{girem la desigualtat}), l'exemple 3 queda

\[
-x^2+4x-3>0 \quad\Rightarrow\quad x^2-4x+3<0,
\]

que té $a>0$. Ara busquem on la paràbola $\cup$ és per sota de l'eix, que és entre les arrels: $(1,3)$. És el mateix resultat.
\end{propietat}

\begin{exemple}{\textbf{Exemple 4:} $x^2>9$ (incompleta)}
La forma canònica és $x^2-9>0$, amb $a=1>0$ ($\cup$). Els punts de tall són

\[
x^2-9=0 \;\Rightarrow\; x^2=9 \;\Rightarrow\; x=\pm 3
\]

\begin{center}
\begin{tikzpicture}
\begin{axis}[grafica, width=0.62\linewidth, height=0.5*0.62\linewidth, xmin=-6, xmax=6, ymin=-11, ymax=8]
\addplot[g-corba, domain=-6:6, samples=160] {1.0*x^2 + 0.0*x + -9.0};
\addplot[g-solcorba, domain=-6:-3, samples=100] {1.0*x^2 + 0.0*x + -9.0};
\draw[g-interval] (axis cs:-6,0) -- (axis cs:-3,0);
\addplot[g-solcorba, domain=3:6, samples=100] {1.0*x^2 + 0.0*x + -9.0};
\draw[g-interval] (axis cs:3,0) -- (axis cs:6,0);
\draw[g-buit] (axis cs:-3,0) circle[radius=2.4pt];
\node[below=3pt, font=\scriptsize, fill=white, inner sep=1pt] at (axis cs:-3,0) {$-3$};
\draw[g-buit] (axis cs:3,0) circle[radius=2.4pt];
\node[below=3pt, font=\scriptsize, fill=white, inner sep=1pt] at (axis cs:3,0) {$3$};
\end{axis}
\end{tikzpicture}
\end{center}

Fem la taula de signes:

\begin{center}
\begin{tabular}{lccccc}
\toprule
$x$ & $(-\infty,-3)$ & $-3$ & $(-3,3)$ & $3$ & $(3,+\infty)$ \\
\midrule
Signe & $+$ & $0$ & $-$ & $0$ & $+$ \\
$>0$? & \textbf{sí} & no & no & no & \textbf{sí} \\
\bottomrule
\end{tabular}
\end{center}

Volem que sigui $>0$: els intervals amb signe $+$, \textbf{fora} de les arrels:

\[
(-\infty,-3)\cup(3,+\infty)
\]

\textbf{Compte:} és un error freqüent escriure «$x>\pm 3$». Per exemple, $x=-4$ és solució, perquè $(-4)^2=16>9$, i en canvi $x=0$ no ho és, perquè $0^2=0\not>9$.
\end{exemple}

\subsection{Cas sense punts de tall}

Quan $\Delta<0$, l'equació $ax^2+bx+c=0$ no té solució i la paràbola \textbf{no talla} l'eix $OX$. Per tant:

\begin{itemize}
\item Si $a>0$ ($\cup$), la paràbola és \textbf{sempre per sobre} de l'eix: $ax^2+bx+c>0$ per a qualsevol $x$.
\item Si $a<0$ ($\cap$), la paràbola és \textbf{sempre per sota} de l'eix: $ax^2+bx+c<0$ per a qualsevol $x$.
\end{itemize}

La inequació, doncs, o bé es compleix per a \textbf{tots} els nombres reals, o bé \textbf{no té cap solució}.

\begin{exemple}{\textbf{Exemple 5:} $x^2+2x+3>0$}
Com que $a=1>0$, la paràbola és \textbf{còncava} ($\cup$). Busquem els punts de tall:

\[
\begin{aligned}
x^2+2x+3&=0\\
\Delta&=2^2-4\cdot 1\cdot 3=4-12=-8<0
\end{aligned}
\]

No hi ha punts de tall: la paràbola és sempre per sobre de l'eix.

\begin{center}
\begin{tikzpicture}
\begin{axis}[grafica, width=0.62\linewidth, height=0.5*0.62\linewidth, xmin=-4, xmax=2, ymin=-2, ymax=8]
\addplot[g-corba, domain=-4:2, samples=160] {1.0*x^2 + 2.0*x + 3.0};
\addplot[g-solcorba, domain=-4:2, samples=100] {1.0*x^2 + 2.0*x + 3.0};
\draw[g-interval] (axis cs:-4,0) -- (axis cs:2,0);
\end{axis}
\end{tikzpicture}
\end{center}

Sense punts de tall, la taula té un sol interval, tota la recta:

\begin{center}
\begin{tabular}{lc}
\toprule
$x$ & $\mathbb{R}$ \\
\midrule
Signe & $+$ \\
$>0$? & \textbf{sí} \\
\bottomrule
\end{tabular}
\end{center}

La desigualtat es compleix sempre. La solució és $\mathbb{R}$, és a dir, $(-\infty,+\infty)$.

En canvi, la inequació $x^2+2x+3\leq 0$ \textbf{no té solució}: la paràbola mai és per sota de l'eix ni el toca.
\end{exemple}

\begin{exemple}{\textbf{Exemple 6:} $-x^2+2x-5\geq 0$}
Ara $a=-1<0$: la paràbola és \textbf{convexa} ($\cap$). Busquem els punts de tall:

\[
\begin{aligned}
-x^2+2x-5&=0\\
\Delta&=2^2-4\cdot(-1)\cdot(-5)\\
&=4-20=-16<0
\end{aligned}
\]

No hi ha punts de tall: la paràbola és sempre per sota de l'eix.

\begin{center}
\begin{tikzpicture}
\begin{axis}[grafica, width=0.62\linewidth, height=0.5*0.62\linewidth, xmin=-2, xmax=4, ymin=-9, ymax=2]
\addplot[g-corba, domain=-2:4, samples=160] {-1.0*x^2 + 2.0*x + -5.0};
\end{axis}
\end{tikzpicture}
\end{center}

La taula té un sol interval:

\begin{center}
\begin{tabular}{lc}
\toprule
$x$ & $\mathbb{R}$ \\
\midrule
Signe & $-$ \\
$\geq 0$? & no \\
\bottomrule
\end{tabular}
\end{center}

Mai és $\geq 0$, i per tant la inequació \textbf{no té solució}. Ho escrivim $\varnothing$ (conjunt buit).

En canvi, $-x^2+2x-5<0$ es compleix per a tots els nombres reals: la solució és $\mathbb{R}$.
\end{exemple}

\subsection{\texorpdfstring{Cas tangent a l'eix $OX$}{Cas tangent a l'eix OX}}

Quan $\Delta=0$, l'equació $ax^2+bx+c=0$ té una sola solució (doble), $x_0$. La paràbola \textbf{toca} l'eix $OX$ en aquest punt, que és el vèrtex, però \textbf{no el travessa}: a banda i banda de $x_0$ queda sempre al mateix costat de l'eix.

\begin{exemple}{\textbf{Exemple 7:} $x^2-4x+4$ comparat amb $0$}
Com que $a=1>0$, la paràbola és \textbf{còncava} ($\cup$). Busquem els punts de tall:

\[
\begin{aligned}
x^2-4x+4&=0\\
\Delta&=(-4)^2-4\cdot 1\cdot 4=16-16=0\\
x&=\frac{4}{2}=2 \quad\text{(arrel doble)}
\end{aligned}
\]

De fet, $x^2-4x+4=(x-2)^2$. La paràbola toca l'eix a $x=2$ i, a la resta de punts, és per sobre. A la taula de signes es veu que el signe \textbf{no canvia} en passar per l'arrel doble:

\begin{center}
\begin{tabular}{lccc}
\toprule
$x$ & $(-\infty,2)$ & $2$ & $(2,+\infty)$ \\
\midrule
Signe & $+$ & $0$ & $+$ \\
\bottomrule
\end{tabular}
\end{center}

\begin{center}
\begin{minipage}[t]{0.48\linewidth}
\centering
\begin{center}
\begin{tikzpicture}
\begin{axis}[grafica, width=0.92\linewidth, height=0.5*0.92\linewidth, xmin=-1, xmax=5, ymin=-2, ymax=6]
\addplot[g-corba, domain=-1:5, samples=160] {1.0*x^2 + -4.0*x + 4.0};
\addplot[g-solcorba, domain=-1:2, samples=100] {1.0*x^2 + -4.0*x + 4.0};
\draw[g-interval] (axis cs:-1,0) -- (axis cs:2,0);
\addplot[g-solcorba, domain=2:5, samples=100] {1.0*x^2 + -4.0*x + 4.0};
\draw[g-interval] (axis cs:2,0) -- (axis cs:5,0);
\draw[g-buit] (axis cs:2,0) circle[radius=2.4pt];
\node[below=3pt, font=\scriptsize, fill=white, inner sep=1pt] at (axis cs:2,0) {$2$};
\end{axis}
\end{tikzpicture}
\end{center}

$x^2-4x+4>0$
\end{minipage}\hfill
\begin{minipage}[t]{0.48\linewidth}
\centering
\begin{center}
\begin{tikzpicture}
\begin{axis}[grafica, width=0.92\linewidth, height=0.5*0.92\linewidth, xmin=-1, xmax=5, ymin=-2, ymax=6]
\addplot[g-corba, domain=-1:5, samples=160] {1.0*x^2 + -4.0*x + 4.0};
\draw[g-ple] (axis cs:2,0) circle[radius=2.4pt];
\node[below=3pt, font=\scriptsize, fill=white, inner sep=1pt] at (axis cs:2,0) {$2$};
\end{axis}
\end{tikzpicture}
\end{center}

$x^2-4x+4\leq 0$
\end{minipage}
\end{center}

Segons la desigualtat, obtenim quatre resultats diferents:

\begin{center}
\begin{tabular}{lll}
\toprule
\textbf{Inequació} & \textbf{Què llegim a la taula?} & \textbf{Solució} \\
\midrule
${x^2-4x+4>0}$ & signe $+$ & ${\mathbb{R}\setminus\{2\}}$, és a dir, $(-\infty,2)\cup(2,+\infty)$ \\
${x^2-4x+4\geq 0}$ & signe $+$ o $0$ & $\mathbb{R}$ \\
${x^2-4x+4<0}$ & cap signe $-$ & no té solució ($\varnothing$) \\
${x^2-4x+4\leq 0}$ & només el $0$, a $x=2$ & $x=2$ \\
\bottomrule
\end{tabular}
\end{center}
\end{exemple}

\begin{exemple}{\textbf{Exemple 8:} $-x^2+6x-9\geq 0$}
Ara $a=-1<0$ ($\cap$) i

\[
\begin{aligned}
-x^2+6x-9&=0\\
x^2-6x+9&=0\\
(x-3)^2&=0 \;\Rightarrow\; x=3
\end{aligned}
\]

La paràbola és per sota de l'eix i només el toca a $x=3$:

\begin{center}
\begin{tabular}{lccc}
\toprule
$x$ & $(-\infty,3)$ & $3$ & $(3,+\infty)$ \\
\midrule
Signe & $-$ & $0$ & $-$ \\
$\geq 0$? & no & \textbf{sí} & no \\
\bottomrule
\end{tabular}
\end{center}

L'únic punt on és $\geq 0$ és el de tangència, i la solució és $x=3$.
\end{exemple}

\begin{nota}{Compte amb l'igual}
En el cas tangent, el punt $x_0$ és solució només si la desigualtat inclou l'igual ($\geq$ o $\leq$). Amb $>$ o $<$ l'hem de treure.
\end{nota}

\subsection{Taula resum}

Escrivim $p(x)=ax^2+bx+c$. Per a $a>0$ (paràbola $\cup$), amb arrels $x_1<x_2$ quan $\Delta>0$ i arrel doble $x_0$ quan $\Delta=0$:

\begin{center}
\begin{tabular}{llll}
\toprule
 & ${\Delta>0}$ & ${\Delta=0}$ & ${\Delta<0}$ \\
\midrule
${p(x)>0}$ & $(-\infty,x_1)\cup(x_2,+\infty)$ & ${\mathbb{R}\setminus\{x_0\}}$ & $\mathbb{R}$ \\
${p(x)\geq 0}$ & $(-\infty,x_1]\cup[x_2,+\infty)$ & $\mathbb{R}$ & $\mathbb{R}$ \\
${p(x)<0}$ & $(x_1,x_2)$ & $\varnothing$ & $\varnothing$ \\
${p(x)\leq 0}$ & $[x_1,x_2]$ & només ${x_0}$ & $\varnothing$ \\
\bottomrule
\end{tabular}
\end{center}

Si $a<0$, podem multiplicar per $-1$ i girar la desigualtat per tornar a aquesta taula. De tota manera, més que memoritzar la taula, és millor fer sempre un \textbf{esbós de la paràbola} i llegir-hi la solució.

\begin{propietat}{I si el grau és més gran?}
Per a inequacions polinòmiques de grau més gran que $2$, factoritzem el polinomi i estudiem el seu signe en cada interval. Ho fem a l'apartat de càlcul de dominis.
\end{propietat}

\section{Nota històrica}

La història de la resolució d'equacions és, en essència, la història de l'àlgebra: l'esforç continu de la humanitat per trobar valors desconeguts a partir de relacions conegudes, un pilar fonamental del pensament científic i tecnològic.

Els primers intents es remunten a les civilitzacions antigues. Els \textbf{babilonis} (cap al 1800 aC) ja resolien problemes pràctics de repartiment de terres o herències que, en termes moderns, equivalen a \textbf{equacions lineals i quadràtiques}, mitjançant receptes algorítmiques, sense la notació simbòlica que fem servir avui. Els grecs, com Euclides i Diofant, van abordar aquests problemes des d'un punt de vista més geomètric.

El pas decisiu cap a l'àlgebra com a disciplina es va donar al \textbf{món arabo-islàmic medieval}. Al segle IX, el matemàtic persa \textbf{Al-Khwarizmí} (d'on prové la paraula "algoritme") va escriure el tractat \emph{Al-Jabr} (d'on prové "àlgebra"), amb mètodes sistemàtics per resoldre equacions lineals i quadràtiques. El seu enfocament, que buscava mètodes generals, va marcar l'inici de l'àlgebra abstracta.

El següent gran salt es va produir a la \textbf{Itàlia del Renaixement} (segle XVI), amb una de les històries més dramàtiques de les matemàtiques: la resolució de les \textbf{equacions cúbiques i quàrtiques}. Matemàtics com \textbf{Tartaglia} i \textbf{Cardano} van trobar la fórmula general per a les equacions de grau 3, i \textbf{Ferrari}, alumne de Cardano, la de grau 4. Irònicament, la fórmula de la cúbica va forçar matemàtics com \textbf{Bombelli} a treballar amb arrels quadrades de nombres negatius, donant el primer impuls al naixement dels nombres complexos.

Durant segles es va buscar una fórmula similar (basada en arrels) per a les equacions de grau 5. El misteri no es va resoldre fins al segle XIX, quan \textbf{Ruffini} i, de manera concloent, \textbf{Niels Henrik Abel}, van demostrar que aquesta fórmula general no existeix. Poc després, \textbf{Évariste Galois} va desenvolupar una teoria que explicava exactament per què algunes equacions polinòmiques tenen solució per radicals i d'altres no, canviant el focus de l'àlgebra de "trobar solucions" a "entendre l'estructura".

Paral·lelament, la necessitat de simplificar càlculs complexos en astronomia i navegació va portar, al segle XVII, \textbf{John Napier} a inventar els \textbf{logaritmes}: un avenç revolucionari, ja que convertia multiplicacions i divisions (feixugues) en sumes i restes. Les \textbf{equacions exponencials} van sorgir com la relació inversa dels logaritmes, i va ser \textbf{Leonhard Euler}, al segle XVIII, qui va consolidar la relació entre les funcions exponencials, logarítmiques i trigonomètriques, donant-los la seva forma moderna.

Avui aquestes equacions són el llenguatge bàsic de la ciència: les polinòmiques descriuen trajectòries; les racionals, relacions i proporcions; les irracionals, distàncies i temps; i les exponencials i logarítmiques modelen tot allò que creix o decreix de forma proporcional a la seva mida, des de la demografia i les finances fins a la desintegració radioactiva o l'escala de pH.
